\studynote{1}{Sun 10 Sep 2023 22:01}{Geometric algebra and Moment of Inertia}

\section{Geometric Algebra and Moment of Inertia}


In this note we present the computational rules in $\mathbb{G}^3$, the three-dimensional geometric algebra, and give explicit examples on how to calculate the moment of inertia in this framework. Moreover, we compare our formulation to the traditional formulation using cross product.

\subsection{The Geometric Algebra}

\begin{definition}
	The 3-dimensional geometric algebra are spanned by the following base elements:
	\begin{itemize}
		\item (Vectors) $e_1, e_2, e_3$
		\item (Bivectors) $e_1 \wedge e_2, e_1 \wedge e_3, e_2 \wedge e_3$
		\item (Trivector / Pseudoscalar) $e_1 \wedge e_2 \wedge e_3 =: i$
	\end{itemize}
	The inner product $\cdot$ and outer product $\wedge$ follow the following rules:
	\begin{enumerate}
		\item Linear.
		\item $\cdot$ is commutative;  $\wedge$ is anti-commutative.
		\item 
	\end{enumerate}
\end{definition}


